\subsection{The surface cases and the unstable Cartesian topoi}

\begin{result}{E.T11}{Known, with full-group and topos reductions proved}{The two remaining stable dimensions}
The field and topos reconstruction statements (E.T19) and (E.T21)
hold for function fields and normal separated varieties over
$k=\overline\Q$ of dimension at least two.  For the two stable
types of modular dimension two, the resulting comparison maps give
\[
 \Aut\bigl((\mathcal M_{0,5,k})_{\mathrm{et}}\bigr)
       =S_5\times G_{\Q},\qquad
 \Aut\bigl((\mathcal M_{1,2,k})_{\mathrm{et}}\bigr)=G_{\Q}.
                                                               \tag{E.T29}
\]
Here all topos equivalences are unringed and taken modulo natural
isomorphism.  The displayed maps are relabelling and semilinear
base change.  In the second formula the relabelling transposition
is naturally isomorphic to the identity algebraic stack functor,
as computed in E.T4.
\end{result}
\begin{proof}[Derivation and source audit]
We use \cite[Theorem 1.1]{OAI2026BogomolovPop}, in the primary
version of 23 September 2026 specified in the bibliography.  Its
input consists of the maximal pro-$\ell$ abelian quotient
$\Pi_F$ and the commutator pairing in the maximal
abelian-by-central quotient $\Pi_F^c$.  In characteristic zero
and relative transcendence degree at least two, its conclusion
is a bijection from field isomorphisms to bracket-compatible
isomorphisms of these abelian quotients, modulo multiplication
by one common unit of $\Z_\ell$.  A continuous full-Galois
isomorphism induces exactly this input: the quotients are
characteristic and its induced maps commute with commutators.
There is no chosen arithmetic action in the hypothesis.

We record the checked steps of the cited proof, since the
two-dimensional range is essential here.  Its local character
construction, Sections 2--3, recovers quasi-prime valuations,
their minimized inertia and decomposition groups, and the
successive residual character spaces.  The alternating-pair
argument was checked against \cite[Section 2 and Theorem 3.3]{Topaz2019Integral}.
It uses characters of multiplicative groups at the residual
steps, so residue characteristic $\ell$ is included at that
stage.  Integral inertia is recovered by intersecting the
rational character lines with the original compact modules.

For a residual curve the local test uses the image, in
$D_v/I_v$, of the intersection of $D_v$ with the closure of
the \emph{union} of the quasi-prime inertia groups.
The exact hypothesis and conclusion of
\cite[Section 4, Theorem 4.2(i)(a)]{Pop2015Minimized}
were opened and checked.  The required conditions are that
$vF$ has no nonzero $\ell$-divisible convex subgroup and that
there is a subfield $k_1$ with
$k_1v=kv$ and
$\operatorname{trdeg}(F/k_1)=\operatorname{trdeg}(Fv/k_1v)=1$.
For the composite quasi-prime valuations in use, the last
relative rank-one step excludes such a convex subgroup.
Choosing elements with rationally independent values modulo
$vk$ gives $k_1$; the unique least monomial of a polynomial
shows $k_1v=kv$.  These are precisely the verifications in
Lemma 3.6 of the 2026 source.

The criterion recognizes a residual curve over an algebraic
closure of a finite field, together with all its point inertia
lines.  We also checked its proof through Lemmas 4.3--4.5 and
Topaz's appended Lemma 5 and Main Theorem.  The latter obtains
the single relation among ordinary point inertia by dualizing
the divisor sequence.  A nontrivial valuation of the constant
field gives a second, disjoint family of residual point inertia,
whose relation is linearly independent of the first.  This
excludes the single-relation condition and proves the needed
converse, including residue characteristic $\ell$.

Sections 4--6 of the 2026 proof were read completely and checked
independently in this work.  The finite root covers preserve
the full kernel of a finite collection of Kummer relations
under specialization.  A complete-intersection curve in one
fixed embedding gives a genus bound independent of the finite
test.  The five-anchor argument first proves actual unit
status on the source curve field.  Its Riemann--Hurwitz
coefficient is $3-5/\ell>0$ even for $\ell=2$; the different
inequality includes wild ramification.  It yields a bound on
the degrees of every exact support divisor, independent of
the Kummer exponent.

The descent to one curve field uses a fixed curve in the
parameter space of those bounded-degree divisors.  For each
finite Kummer test, normalization in a finite splitting
extension of the parameter function field makes its root cover
generically split.  On the relative smooth locus the base
change is normal, and a generically split finite etale cover
is split everywhere.  Consequently the generic point of the
already fixed parameter curve can be used for every test.
This is the step in Lemma 6.3 that supplies one curve field
for the entire completed subgroup.  The finite-extension
norm and saturation then descend it to the original field.
These arguments give the correspondence of curve subfields,
true divisorial inertia, and rational quotients in
Theorem 6.1 and Propositions 6.6--6.7.

For the final step we opened the primary author versions of
\cite[Theorem 2.1(2)]{Pop2012BogomolovI} and
\cite[Fact/Definition 43 and Theorem 45]{Pop2012DecompositionGraphs},
including the latter's proof in Section 6.  The global Isom
statement assumes transcendence degree greater than one over
an algebraically closed field and compatibility with a
Bertini family of rational quotient diagrams.  Its constant
field has no finite-field restriction.  All rational quotients
recovered above contain the required Bertini family; their
restriction maps commute because they are dual to the actual
completed multiplicative-subfield inclusions.  Pop's proof
recovers the multiplicative lattice and its projective lines,
then the field by projective geometry.  Thus every hypothesis
of the global statement is supplied in relative dimension two.
The source locations and the independent geometric audit are
retained in the work records.  Silberstein's meridian assertion,
discussed in E.T3, is not an input to this derivation.

This proves that the smaller-quotient reconstruction used at
the start of E.T8 is available in transcendence degree two as
well.  Every finite extension of such a field has the same
transcendence degree.  Hence the construction there applies to
each open normal subgroup, with the same finite-index and
torsion-free comparison.  Its compatible field isomorphisms
assemble over the separable closures and intertwine the
original full Galois isomorphism.  This proves (E.T19) in the
claimed range, with its stated uniqueness and all open-subgroup
compatibilities.

The proof of E.T9 uses this birational result on the generic
field of each normal separated etale object; all those fields
still have dimension two.  Its inertia-stalk comparison,
intersection of height-one valuation rings, and generating-site
argument therefore give (E.T21) with that same range.  Finally
E.T10 applies to the finite level-$3$ atlas and its iterated
etale fibre products for $\mathcal M_{1,2}$; their dimensions
are two.  For $\mathcal M_{0,5}$ the atlas is the scheme itself.
Descent gives all semilinear algebraic stack automorphisms.
E.T4 gives their $k$-linear groups $1$ and $S_5$, respectively.
Their models over $\Q$ split the base-field map, and relabelling
is defined over $\Q$, proving the direct products in (E.T29).
\end{proof}

We now specify a sheaf convention for the four unstable types.
For an algebraic stack $\mathcal X$, let
$\operatorname{CartEt}(\mathcal X)$ be the category of sheaves
of sets on a smooth presentation, with the inverse-image
descent isomorphism on its relation space and the cocycle
condition on its triple fibre product.  Equivalently these are
Cartesian sheaves for that presentation's simplicial etale
sites.  The inverse constructions are given in
\cite[Lemma 85.12.4]{StacksCartesianDescent}.
Effective smooth descent for etale sheaves identifies the
categories for two presentations, by their common refinement.
For a Deligne--Mumford stack we may choose an etale presentation;
ordinary etale descent then identifies this category with its
full small etale topos, by the restriction and descent functors.
Thus the convention agrees with every stable computation above.

\begin{result}{E.T12}{Proved}{Connected bands and the four unstable types}
Let $S$ be an algebraic stack with a smooth scheme presentation
and let $H\to S$ be a smooth group scheme, of finite presentation,
with geometrically connected fibres.  The projection and its
trivial-torsor section induce inverse equivalences
\[
 \operatorname{CartEt}(S)
       \ \simeq\ \operatorname{CartEt}(B_S H).                \tag{E.T30}
\]
For $k=\overline\Q$ this gives the following complete unstable
calculations and the indicated algebraic comparison maps:
\[
\begin{array}{c|c|c}
 (g,n)&\operatorname{CartEt}(\mathcal M_{g,n,k})
 &\Aut\operatorname{CartEt}(\mathcal M_{g,n,k})
 \\ \hline
 (0,0)&\mathbf{Set}&1\\
 (0,1)&\mathbf{Set}&1\\
 (0,2)&\mathbf{Set}&1\\
 (1,0)&(\mathcal M_{1,1,k})_{\mathrm{et}}&C_2\times\mathcal A.
\end{array}                                                   \tag{E.T31}
\]
The group $\mathcal A$ and its exact inertia conditions are
those of E.T7.  In the last row the equivalence is pullback
along the Jacobian map.  All four comparison homomorphisms
from the $k$-linear algebraic automorphism groups are trivial.
Their domains, in the order displayed, are $1,1,C_2,C_2$.
\end{result}
\begin{proof}
We first work over a scheme $S$.  Write $p:H\to S$ for the
structure map.  Its flatness, quasi-compactness and geometric
connectedness imply
\[
          \mathcal F\xrightarrow{\sim}p_*p^*\mathcal F
                 \quad(\mathcal F\in S_{\rm et}).            \tag{E.T32}
\]
Indeed, evaluate on an arbitrary etale $U\to S$ and apply
\cite[Lemma 59.39.3]{StacksConnectedPullback}
to $H\times_S U\to U$ and $\mathcal F|_U$.
The lemma's descent proof applies to every sheaf of sets:
on each connected geometric fibre a section of the pulled-back
constant sheaf is constant; the two pullbacks to the fibre
square agree and effective fpqc descent gives the section on
$U$.  Thus (E.T32) holds for all etale $U$, and $p^*$ is fully
faithful by adjunction.

The smooth presentation $S\to B_S H$ has relation space $H$,
with both structure maps equal to $p$.  A Cartesian sheaf is
therefore a pair $(\mathcal F,\alpha)$ with
$\alpha:p^*\mathcal F\to p^*\mathcal F$ an isomorphism
satisfying the group descent cocycle.  Full faithfulness gives
a unique $\beta:\mathcal F\to\mathcal F$ with
$\alpha=p^*\beta$.  Pullback along the identity section makes
$\beta$ the identity: the cocycle there says
$\beta=\beta\circ\beta$, and $\beta$ is invertible.
It follows that the only descent datum is the trivial one.
Morphisms descend uniquely by the same full faithfulness.
Projection pullback and restriction to the trivial torsor are
therefore quasi-inverse.  For an algebraic stack $S$ apply
this construction on a smooth scheme atlas and all its fibre
products.  All equivalences just constructed are pullbacks
along the stated maps, so they commute with that descent
and give (E.T30).

E.T5 identifies the first three moduli stacks with
$BPGL_2$, $B(\mathbf G_a\rtimes\mathbf G_m)$, and
$B\mathbf G_m$, respectively.  Their displayed bands are
smooth and connected, so (E.T30) identifies their Cartesian
sheaves with $\operatorname{Spec}(k)_{\rm et}=\mathbf{Set}$.
Every equivalence of $\mathbf{Set}$ is naturally isomorphic
to the identity: it preserves the terminal singleton and all
its coproducts, and every set is such a coproduct.
This proves the first three groups and their comparison maps.

For the last row, E.T6 gives the actual equivalence
$\mathcal M_{1,0}=B_{\mathcal M_{1,1}}\mathcal E$.
The universal elliptic curve $\mathcal E$ is smooth, proper,
and geometrically connected.  Equation (E.T30) consequently
identifies its Cartesian sheaves with the small etale topos
of $\mathcal M_{1,1}$.  E.T7 computes the resulting
automorphism group.  The nontrivial algebraic automorphism
of E.T6 tensors an elliptic torsor with the specified universal
torsor and induces the identity on the Jacobian base.  Under
(E.T30) its inverse-image functor is consequently naturally
the identity, proving the last comparison map.  In particular,
this comparison records the entire algebraic kernel.
\end{proof}
